\documentclass{article}

\usepackage[margin=1in]{geometry}
\usepackage{hyperref}
\usepackage{xcolor}
\usepackage{listings}
\usepackage{longtable}
\usepackage{booktabs}
\usepackage{enumitem}
\usepackage{titlesec}
\usepackage{fancyhdr}

\hypersetup{
  colorlinks=true,
  linkcolor=blue,
  citecolor=blue,
  urlcolor=blue
}

\pagestyle{fancy}
\fancyhf{}
\fancyhead[L]{RFC 8829}
\fancyhead[R]{JSEP}
\fancyfoot[C]{\thepage}

\title{JavaScript Session Establishment Protocol (JSEP)\\
\large RFC 8829\\
Internet Engineering Task Force (IETF)\\
Category: Standards Track}

\author{
J. Uberti (Google) \and
C. Jennings (Cisco) \and
E. Rescorla, Ed. (Mozilla)
}

\date{January 2021}

\begin{document}

\maketitle

\begin{abstract}
This document describes the mechanisms for allowing a JavaScript
application to control the signaling plane of a multimedia session
via the interface specified in the W3C RTCPeerConnection API and
discusses how this relates to existing signaling protocols.
\end{abstract}

\subsection*{Status of This Memo}

This is an Internet Standards Track document.

This document is a product of the Internet Engineering Task Force
(IETF).  It represents the consensus of the IETF community.  It has
received public review and has been approved for publication by the
Internet Engineering Steering Group (IESG).  Further information on
Internet Standards is available in Section 2 of RFC 7841.

Information about the current status of this document, any errata,
and how to provide feedback on it may be obtained at
\url{https://www.rfc-editor.org/info/rfc8829}.

\subsection*{Copyright Notice}

Copyright (c) 2021 IETF Trust and the persons identified as the
document authors.  All rights reserved.

This document is subject to BCP 78 and the IETF Trust's Legal
Provisions Relating to IETF Documents
(\url{https://trustee.ietf.org/license-info}) in effect on the date of
publication of this document.  Please review these documents
carefully, as they describe your rights and restrictions with respect
to this document.  Code Components extracted from this document must
include Simplified BSD License text as described in Section 4.e of
the Trust Legal Provisions and are provided without warranty as
described in the Simplified BSD License.

\tableofcontents
\newpage

%% =============================================================
\section{Introduction}

This document describes how the W3C Web Real-Time Communication
(WebRTC) RTCPeerConnection interface [W3C.webrtc] is used to control
the setup, management, and teardown of a multimedia session.

\subsection{General Design of JSEP}

WebRTC call setup has been designed to focus on controlling the media
plane, leaving signaling-plane behavior up to the application as much
as possible.  The rationale is that different applications may prefer
to use different protocols, such as the existing SIP call signaling
protocol, or something custom to the particular application, perhaps
for a novel use case.  In this approach, the key information that
needs to be exchanged is the multimedia session description, which
specifies the transport and media configuration information necessary
to establish the media plane.

With these considerations in mind, this document describes the
JavaScript Session Establishment Protocol (JSEP), which allows for
full control of the signaling state machine from JavaScript.  As
described above, JSEP assumes a model in which a JavaScript
application executes inside a runtime containing WebRTC APIs (the
``JSEP implementation'').  The JSEP implementation is almost entirely
divorced from the core signaling flow, which is instead handled by
the JavaScript making use of two interfaces: (1) passing in local and
remote session descriptions and (2) interacting with the Interactive
Connectivity Establishment (ICE) state machine [RFC8445].  The
combination of the JSEP implementation and the JavaScript application
is referred to throughout this document as a ``JSEP endpoint''.

In this document, the use of JSEP is described as if it always occurs
between two JSEP endpoints.  Note, though, that in many cases it will
actually be between a JSEP endpoint and some kind of server, such as
a gateway or Multipoint Control Unit (MCU).  This distinction is
invisible to the JSEP endpoint; it just follows the instructions it
is given via the API.

JSEP's handling of session descriptions is simple and
straightforward.  Whenever an offer/answer exchange is needed, the
initiating side creates an offer by calling a createOffer API.  The
application then uses that offer to set up its local configuration
via the setLocalDescription API.  The offer is finally sent off to
the remote side over its preferred signaling mechanism (e.g.,
WebSockets); upon receipt of that offer, the remote party installs it
using the setRemoteDescription API.

To complete the offer/answer exchange, the remote party uses the
createAnswer API to generate an appropriate answer, applies it using
the setLocalDescription API, and sends the answer back to the
initiator over the signaling channel.  When the initiator gets that
answer, it installs it using the setRemoteDescription API, and
initial setup is complete.  This process can be repeated for
additional offer/answer exchanges.

Regarding ICE [RFC8445], JSEP decouples the ICE state machine from
the overall signaling state machine.  The ICE state machine must
remain in the JSEP implementation because only the implementation has
the necessary knowledge of candidates and other transport
information.  Performing this separation provides additional
flexibility in protocols that decouple session descriptions from
transport.  For instance, in traditional SIP, each offer or answer is
self-contained, including both the session descriptions and the
transport information.  However, [RFC8840] allows SIP to be used with
Trickle ICE [RFC8838], in which the session description can be sent
immediately and the transport information can be sent when available.
Sending transport information separately can allow for faster ICE and
DTLS startup, since ICE checks can start as soon as any transport
information is available rather than waiting for all of it.  JSEP's
decoupling of the ICE and signaling state machines allows it to
accommodate either model.

Although it abstracts signaling, the JSEP approach requires that the
application be aware of the signaling process.  While the application
does not need to understand the contents of session descriptions to
set up a call, the application must call the right APIs at the right
times, convert the session descriptions and ICE information into the
defined messages of its chosen signaling protocol, and perform the
reverse conversion on the messages it receives from the other side.

One way to make life easier for the application is to provide a
JavaScript library that hides this complexity from the developer;
said library would implement a given signaling protocol along with
its state machine and serialization code, presenting a higher-level
call-oriented interface to the application developer.  For example,
libraries exist to provide implementations of the SIP [RFC3261] and
Extensible Messaging and Presence Protocol (XMPP) [RFC6120] signaling
protocols atop the JSEP API.  Thus, JSEP provides greater control for
the experienced developer without forcing any additional complexity
on the novice developer.

\subsection{Other Approaches Considered}

One approach that was considered instead of JSEP was to include a
lightweight signaling protocol.  Instead of providing session
descriptions to the API, the API would produce and consume messages
from this protocol.  While providing a more high-level API, this put
more control of signaling within the JSEP implementation, forcing it
to have to understand and handle concepts like signaling glare (see
[RFC3264], Section 4).

A second approach that was considered but not chosen was to decouple
the management of the media control objects from session
descriptions, instead offering APIs that would control each component
directly.  This was rejected based on the argument that requiring
exposure of this level of complexity to the application programmer
would not be beneficial; it would (1) result in an API where even a
simple example would require a significant amount of code to
orchestrate all the needed interactions and (2) create a large API
surface that would need to be agreed upon and documented.  In
addition, these API points could be called in any order, resulting in
a more complex set of interactions with the media subsystem than the
JSEP approach, which specifies how session descriptions are to be
evaluated and applied.

One variation on JSEP that was considered was to keep the basic
session-description-oriented API but to move the mechanism for
generating offers and answers out of the JSEP implementation.
Instead of providing createOffer/createAnswer methods within the
implementation, this approach would instead expose a getCapabilities
API, which would provide the application with the information it
needed in order to generate its own session descriptions.  This
increases the amount of work that the application needs to do; it
needs to know how to generate session descriptions from capabilities,
and especially how to generate the correct answer from an arbitrary
offer and the supported capabilities.  While this could certainly be
addressed by using a library like the one mentioned above, it
basically forces the use of said library even for a simple example.
Providing createOffer/createAnswer avoids this problem.

\subsection{Contradiction regarding bundle-only ``m='' sections}

Since the approval of the WebRTC specification documents, the IETF
has become aware of an inconsistency between the document specifying
JSEP and the document specifying BUNDLE (this RFC and [RFC8843],
respectively).  Rather than delaying publication further to come to a
resolution, the documents are being published as they were originally
approved.  The IETF intends to restart work on these technologies,
and revised versions of these documents will be published as soon as
a resolution becomes available.

The specific issue involves the handling of ``m='' sections that are
designated as bundle-only, as discussed in Section 4.1.1 of this RFC.
Currently, there is divergence between JSEP and BUNDLE, as well as
between these specifications and existing browser implementations:

\begin{itemize}
  \item JSEP prescribes that said ``m='' sections should use port zero and
    add an ``a=bundle-only'' attribute in initial offers, but not in
    answers or subsequent offers.

  \item BUNDLE prescribes that these ``m='' sections should be marked as
    described in the previous point, but in all offers and answers.

  \item Most current browsers do not mark any ``m='' sections with port zero
    and instead use the same port for all bundled ``m='' sections; some
    others follow the JSEP behavior.
\end{itemize}

%% =============================================================
\section{Terminology}

The key words ``MUST'', ``MUST NOT'', ``REQUIRED'', ``SHALL'', ``SHALL NOT'',
``SHOULD'', ``SHOULD NOT'', ``RECOMMENDED'', ``NOT RECOMMENDED'', ``MAY'', and
``OPTIONAL'' in this document are to be interpreted as described in
BCP~14 [RFC2119] [RFC8174] when, and only when, they appear in all
capitals, as shown here.

%% =============================================================
\section{Semantics and Syntax}

\subsection{Signaling Model}

JSEP does not specify a particular signaling model or state machine,
other than the generic need to exchange session descriptions in the
fashion described by [RFC3264] (offer/answer) in order for both sides
of the session to know how to conduct the session.  JSEP provides
mechanisms to create offers and answers, as well as to apply them to
a session.  However, the JSEP implementation is totally decoupled
from the actual mechanism by which these offers and answers are
communicated to the remote side, including addressing,
retransmission, forking, and glare handling.  These issues are left
entirely up to the application; the application has complete control
over which offers and answers get handed to the implementation, and
when.

\begin{verbatim}
        +-----------+                               +-----------+
        |  Web App  |<--- App-Specific Signaling -->|  Web App  |
        +-----------+                               +-----------+
              ^                                            ^
              |  SDP                                       |  SDP
              V                                            V
        +-----------+                                +-----------+
        |   JSEP    |<----------- Media ------------>|   JSEP    |
        |   Impl.   |                                |   Impl.   |
        +-----------+                                +-----------+

                      Figure 1: JSEP Signaling Model
\end{verbatim}

\subsection{Session Descriptions and State Machine}

In order to establish the media plane, the JSEP implementation needs
specific parameters to indicate what to transmit to the remote side,
as well as how to handle the media that is received.  These
parameters are determined by the exchange of session descriptions in
offers and answers, and there are certain details to this process
that must be handled in the JSEP APIs.

Whether a session description applies to the local side or the remote
side affects the meaning of that description.  For example, the list
of codecs sent to a remote party indicates what the local side is
willing to receive, which, when intersected with the set of codecs
the remote side supports, specifies what the remote side should send.
However, not all parameters follow this rule; some parameters are
declarative, and the remote side must either accept them or reject
them altogether.  An example of such a parameter is the TLS
fingerprints [RFC8122] as used in the context of DTLS [RFC6347];
these fingerprints are calculated based on the local certificate(s)
offered and are not subject to negotiation.

In addition, various RFCs put different conditions on the format of
offers versus answers.  For example, an offer may propose an
arbitrary number of ``m='' sections (i.e., media descriptions as
described in [RFC4566], Section 5.14), but an answer must contain the
exact same number as the offer.

Lastly, while the exact media parameters are known only after an
offer and an answer have been exchanged, the offerer may receive ICE
checks, and possibly media (e.g., in the case of a re-offer after a
connection has been established) before it receives an answer.  To
properly process incoming media in this case, the offerer's media
handler must be aware of the details of the offer before the answer
arrives.

Therefore, in order to handle session descriptions properly, the JSEP
implementation needs:

\begin{enumerate}
  \item To know if a session description pertains to the local or remote
    side.

  \item To know if a session description is an offer or an answer.

  \item To allow the offer to be specified independently of the answer.
\end{enumerate}

JSEP addresses this by adding both setLocalDescription and
setRemoteDescription methods and having session description objects
contain a type field indicating the type of session description being
supplied.  This satisfies the requirements listed above for both the
offerer, who first calls setLocalDescription(sdp [offer]) and then
later setRemoteDescription(sdp [answer]), and the answerer, who first
calls setRemoteDescription(sdp [offer]) and then later
setLocalDescription(sdp [answer]).

During the offer/answer exchange, the outstanding offer is considered
to be ``pending'' at the offerer and the answerer, as it may be either
accepted or rejected.  If this is a re-offer, each side will also
have ``current'' local and remote descriptions, which reflect the
result of the last offer/answer exchange.  Sections 4.1.14, 4.1.16,
4.1.13, and 4.1.15 provide more detail on pending and current
descriptions.

JSEP also allows for an answer to be treated as provisional by the
application.  Provisional answers provide a way for an answerer to
communicate initial session parameters back to the offerer, in order
to allow the session to begin, while allowing a final answer to be
specified later.  This concept of a final answer is important to the
offer/answer model; when such an answer is received, any extra
resources allocated by the caller can be released, now that the exact
session configuration is known.  These ``resources'' can include things
like extra ICE components, Traversal Using Relays around NAT (TURN)
candidates, or video decoders.  Provisional answers, on the other
hand, do no such deallocation; as a result, multiple dissimilar
provisional answers, with their own codec choices, transport
parameters, etc., can be received and applied during call setup.
Note that the final answer itself may be different than any received
provisional answers.

In [RFC3264], the constraint at the signaling level is that only one
offer can be outstanding for a given session, but at the JSEP level,
a new offer can be generated at any point.  For example, when using
SIP for signaling, if one offer is sent and is then canceled using a
SIP CANCEL, another offer can be generated even though no answer was
received for the first offer.  To support this, the JSEP media layer
can provide an offer via the createOffer method whenever the
JavaScript application needs one for the signaling.  The answerer can
send back zero or more provisional answers and then finally end the
offer/answer exchange by sending a final answer.  The state machine
for this is as follows:

\begin{verbatim}
                      setRemote(OFFER)               setLocal(PRANSWER)
                          /-----\                               /-----\
                          |     |                               |     |
                          v     |                               v     |
           +---------------+    |                +---------------+    |
           |               |----/                |               |----/
           |  have-        | setLocal(PRANSWER)  | have-         |
           |  remote-offer |------------------- >| local-pranswer|
           |               |                     |               |
           |               |                     |               |
           +---------------+                     +---------------+
                ^   |                                   |
                |   | setLocal(ANSWER)                  |
  setRemote(OFFER)  |                                   |
                |   V                  setLocal(ANSWER) |
           +---------------+                            |
           |               |                            |
           |               |<---------------------------+
           |    stable     |
           |               |<---------------------------+
           |               |                            |
           +---------------+          setRemote(ANSWER) |
                ^   |                                   |
                |   | setLocal(OFFER)                   |
  setRemote(ANSWER) |                                   |
                |   V                                   |
           +---------------+                     +---------------+
           |               |                     |               |
           |  have-        | setRemote(PRANSWER) |have-          |
           |  local-offer  |------------------- >|remote-pranswer|
           |               |                     |               |
           |               |----\                |               |----\
           +---------------+    |                +---------------+    |
                          ^     |                               ^     |
                          |     |                               |     |
                          \-----/                               \-----/
                      setLocal(OFFER)               setRemote(PRANSWER)

                       Figure 2: JSEP State Machine
\end{verbatim}

Aside from these state transitions, there is no other difference
between the handling of provisional (``pranswer'') and final (``answer'')
answers.

\subsection{Session Description Format}

JSEP's session descriptions use Session Description Protocol (SDP)
syntax for their internal representation.  While this format is not
optimal for manipulation from JavaScript, it is widely accepted and
is frequently updated with new features; any alternate encoding of
session descriptions would have to keep pace with the changes to SDP,
at least until the time that this new encoding eclipsed SDP in
popularity.

However, to provide for future flexibility, the SDP syntax is
encapsulated within a SessionDescription object, which can be
constructed from SDP and be serialized out to SDP.  If future
specifications agree on a JSON format for session descriptions, we
could easily enable this object to generate and consume that JSON.

As detailed below, most applications should be able to treat the
SessionDescriptions produced and consumed by these various API calls
as opaque blobs; that is, the application will not need to read or
change them.

\subsection{Session Description Control}

In order to give the application control over various common session
parameters, JSEP provides control surfaces that tell the JSEP
implementation how to generate session descriptions.  This avoids the
need for JavaScript to modify session descriptions in most cases.

Changes to these objects result in changes to the session
descriptions generated by subsequent createOffer/createAnswer calls.

\subsubsection{RtpTransceivers}

RtpTransceivers allow the application to control the RTP media
associated with one ``m='' section.  Each RtpTransceiver has an
RtpSender and an RtpReceiver, which an application can use to control
the sending and receiving of RTP media.  The application may also
modify the RtpTransceiver directly, for instance, by stopping it.

RtpTransceivers generally have a 1:1 mapping with ``m='' sections,
although there may be more RtpTransceivers than ``m='' sections when
RtpTransceivers are created but not yet associated with an ``m=''
section, or if RtpTransceivers have been stopped and disassociated
from ``m='' sections.  An RtpTransceiver is said to be associated with
an ``m='' section if its media identification (mid) property is non-null;
otherwise, it is said to be disassociated.  The associated ``m=''
section is determined using a mapping between transceivers and ``m=''
section indices, formed when creating an offer or applying a remote
offer.

An RtpTransceiver is never associated with more than one ``m=''
section, and once a session description is applied, an ``m='' section
is always associated with exactly one RtpTransceiver.  However, in
certain cases where an ``m='' section has been rejected, as discussed
in Section 5.2.2 below, that ``m='' section will be ``recycled'' and
associated with a new RtpTransceiver with a new MID value.

RtpTransceivers can be created explicitly by the application or
implicitly by calling setRemoteDescription with an offer that adds
new ``m='' sections.

\subsubsection{RtpSenders}

RtpSenders allow the application to control how RTP media is sent.
An RtpSender is conceptually responsible for the outgoing RTP
stream(s) described by an ``m='' section.  This includes encoding the
attached MediaStreamTrack, sending RTP media packets, and generating/
processing the RTP Control Protocol (RTCP) for the outgoing RTP
streams(s).

\subsubsection{RtpReceivers}

RtpReceivers allow the application to inspect how RTP media is
received.  An RtpReceiver is conceptually responsible for the
incoming RTP stream(s) described by an ``m='' section.  This includes
processing received RTP media packets, decoding the incoming
stream(s) to produce a remote MediaStreamTrack, and generating/
processing RTCP for the incoming RTP stream(s).

\subsection{ICE}

\subsubsection{ICE Gathering Overview}

JSEP gathers ICE candidates as needed by the application.  Collection
of ICE candidates is referred to as a gathering phase, and this is
triggered either by the addition of a new or recycled ``m='' section to
the local session description or by new ICE credentials in the
description, indicating an ICE restart.  Use of new ICE credentials
can be triggered explicitly by the application or implicitly by the
JSEP implementation in response to changes in the ICE configuration.

When the ICE configuration changes in a way that requires a new
gathering phase, a `needs-ice-restart' bit is set.  When this bit is
set, calls to the createOffer API will generate new ICE credentials.
This bit is cleared by a call to the setLocalDescription API with new
ICE credentials from either an offer or an answer, i.e., from either
a locally or remotely initiated ICE restart.

When a new gathering phase starts, the ICE agent will notify the
application that gathering is occurring through a state change event.
Then, when each new ICE candidate becomes available, the ICE agent
will supply it to the application via an onicecandidate event; these
candidates will also automatically be added to the current and/or
pending local session description.  Finally, when all candidates have
been gathered, a final onicecandidate event will be dispatched to
signal that the gathering process is complete.

Note that gathering phases only gather the candidates needed by
new/recycled/restarting ``m='' sections; other ``m='' sections continue
to use their existing candidates.  Also, if an ``m='' section is
bundled (either by a successful bundle negotiation or by being marked
as bundle-only), then candidates will be gathered and exchanged for
that ``m='' section if and only if its MID item is a BUNDLE-tag, as
described in [RFC8843].

\subsubsection{ICE Candidate Trickling}

Candidate trickling is a technique through which a caller may
incrementally provide candidates to the callee after the initial
offer has been dispatched; the semantics of ``Trickle ICE'' are defined
in [RFC8838].  This process allows the callee to begin acting upon
the call and setting up the ICE (and perhaps DTLS) connections
immediately, without having to wait for the caller to gather all
possible candidates.  This results in faster media setup in cases
where gathering is not performed prior to initiating the call.

JSEP supports optional candidate trickling by providing APIs, as
described above, that provide control and feedback on the ICE
candidate gathering process.  Applications that support candidate
trickling can send the initial offer immediately and send individual
candidates when they get notified of a new candidate; applications
that do not support this feature can simply wait for the indication
that gathering is complete, and then create and send their offer,
with all the candidates, at that time.

Upon receipt of trickled candidates, the receiving application will
supply them to its ICE agent.  This triggers the ICE agent to start
using the new remote candidates for connectivity checks.

\paragraph{ICE Candidate Format}

In JSEP, ICE candidates are abstracted by an IceCandidate object, and
as with session descriptions, SDP syntax is used for the internal
representation.

The candidate details are specified in an IceCandidate field, using
the same SDP syntax as the ``candidate-attribute'' field defined in
[RFC8839], Section 5.1.  Note that this field does not contain an
``a='' prefix, as indicated in the following example:

\begin{verbatim}
candidate:1 1 UDP 1694498815 192.0.2.33 10000 typ host
\end{verbatim}

The IceCandidate object contains a field to indicate which ICE
username fragment (ufrag) it is associated with, as defined in
[RFC8839], Section 5.4.  This value is used to determine which
session description (and thereby which gathering phase) this
IceCandidate belongs to, which helps resolve ambiguities during ICE
restarts.  If this field is absent in a received IceCandidate
(perhaps when communicating with a non-JSEP endpoint), the most
recently received session description is assumed.

The IceCandidate object also contains fields to indicate which ``m=''
section it is associated with, which can be identified in one of two
ways: either by an ``m='' section index or by a MID.  The ``m='' section
index is a zero-based index, with index N referring to the N+1th ``m=''
section in the session description referenced by this IceCandidate.
The MID is a ``media stream identification'' value, as defined in
[RFC5888], Section 4, which provides a more robust way to identify
the ``m='' section in the session description, using the MID of the
associated RtpTransceiver object (which may have been locally
generated by the answerer when interacting with a non-JSEP endpoint
that does not support the MID attribute, as discussed in Section 5.10
below).  If the MID field is present in a received IceCandidate, it
MUST be used for identification; otherwise, the ``m='' section index is
used instead.

Implementations MUST be prepared to receive objects with some fields
missing, as mentioned above.

\subsubsection{ICE Candidate Policy}

Typically, when gathering ICE candidates, the JSEP implementation
will gather all possible forms of initial candidates --- host, server-reflexive,
and relay.  However, in certain cases, applications may
want to have more specific control over the gathering process, due to
privacy or related concerns.  For example, one may want to only use
relay candidates, to leak as little location information as possible
(keeping in mind that this choice comes with corresponding
operational costs).  To accomplish this, JSEP allows the application
to restrict which ICE candidates are used in a session.  Note that
this filtering is applied on top of any restrictions the
implementation chooses to enforce regarding which IP addresses are
permitted for the application, as discussed in [RFC8828].

There may also be cases where the application wants to change which
types of candidates are used while the session is active.  A prime
example is where a callee may initially want to use only relay
candidates, to avoid leaking location information to an arbitrary
caller, but then change to use all candidates (for lower operational
cost) once the user has indicated that they want to take the call.
For this scenario, the JSEP implementation MUST allow the candidate
policy to be changed in mid-session, subject to the aforementioned
interactions with local policy.

To administer the ICE candidate policy, the JSEP implementation will
determine the current setting at the start of each gathering phase.
Then, during the gathering phase, the implementation MUST NOT expose
candidates disallowed by the current policy to the application, use
them as the source of connectivity checks, or indirectly expose them
via other fields, such as the raddr/rport attributes for other ICE
candidates.  Later, if a different policy is specified by the
application, the application can apply it by kicking off a new
gathering phase via an ICE restart.

\subsubsection{ICE Candidate Pool}

JSEP applications typically inform the JSEP implementation to begin
ICE gathering via the information supplied to setLocalDescription, as
the local description indicates the number of ICE components that
will be needed and for which candidates must be gathered.  However,
to accelerate cases where the application knows the number of ICE
components to use ahead of time, it may ask the implementation to
gather a pool of potential ICE candidates to help ensure rapid media
setup.

When setLocalDescription is eventually called and the JSEP
implementation prepares to gather the needed ICE candidates, it
SHOULD start by checking if any candidates are available in the pool.
If there are candidates in the pool, they SHOULD be handed to the
application immediately via the ICE candidate event.  If the pool
becomes depleted, either because a larger-than-expected number of ICE
components are used or because the pool has not had enough time to
gather candidates, the remaining candidates are gathered as usual.
This only occurs for the first offer/answer exchange, after which the
candidate pool is emptied and no longer used.

One example of where this concept is useful is an application that
expects an incoming call at some point in the future, and wants to
minimize the time it takes to establish connectivity, to avoid
clipping of initial media.  By pre-gathering candidates into the
pool, it can exchange and start sending connectivity checks from
these candidates almost immediately upon receipt of a call.  Note,
though, that by holding on to these pre-gathered candidates, which
will be kept alive as long as they may be needed, the application
will consume resources on the STUN/TURN servers it is using.  (``STUN''
stands for ``Session Traversal Utilities for NAT''.)

\subsubsection{ICE Versions}

While this specification formally relies on [RFC8445], at the time of
its publication, the majority of WebRTC implementations support the
version of ICE described in [RFC5245].  The ``ice2'' attribute defined
in [RFC8445] can be used to detect the version in use by a remote
endpoint and to provide a smooth transition from the older
specification to the newer one.  Implementations MUST be able to
accept remote descriptions that do not have the ``ice2'' attribute.

